% ECONOMICS_PROBLEM: {"id": "M37-645", "title": "Incremental effects of digital advertising", "jel_code": "M37", "source_id": "OP-0075"}
% !TeX program = xelatex
\documentclass[11pt,a4paper]{article}
\usepackage[margin=23mm]{geometry}
\usepackage{fontspec}
\setmainfont{Latin Modern Roman}
\usepackage{amsmath,amssymb,mathtools}
\usepackage{xcolor,enumitem,booktabs,longtable,array,xurl,hyperref}
\definecolor{muted}{HTML}{53636C}
\hypersetup{colorlinks=true,urlcolor=blue,linkcolor=blue}
\setlength{\parindent}{0pt}
\setlength{\parskip}{5pt}
\newcommand{\R}{\mathbb R}
\newcommand{\E}{\mathbb E}
\newcommand{\N}{\mathbb N}
\newcommand{\ind}{\mathbf 1}
\newcommand{\Var}{\operatorname{Var}}
\newcommand{\Cov}{\operatorname{Cov}}
\newcommand{\supp}{\operatorname{supp}}
\DeclareMathOperator*{\argmin}{arg\,min}
\DeclareMathOperator*{\argmax}{arg\,max}
\newcommand{\entrymeta}[3]{{\sffamily\small\color{muted}\textbf{\texttt{#1}}\quad|\quad #2\par #3\par}}
\newcommand{\Problemref}[1]{\texttt{#1}}
\newcommand{\JELref}[2]{\texttt{#2} (JEL \texttt{#1})}
\begin{document}
\section*{Incremental effects of digital advertising}
\textbf{Problem ID:} \texttt{M37-645}\quad \textbf{JEL:} \texttt{M37}\par
% BEGIN ECONOMICS STATEMENT
\label{problem:OP-0075}
\label{atlas:prob:I5}

\noindent\textbf{Original identifier:} I5 (atlas item 75). \textbf{Classification:} Advertising causal measurement research program. \textbf{Original tractability label:} Easy (editorial, not a complexity result).

\subsection*{Definitions and assumptions}

For a baseline eligible consumer population, let $Z\in\{0,1\}$ indicate randomized eligibility for a specified advertising campaign. Let $Y(z)$ be purchases or contribution margin over a prespecified horizon, including all measured channels. Ad exposure $A$ is generally a post-assignment variable determined by auctions, browsing, and competing campaigns. The intention-to-treat estimand is $\tau_Z=\mathbb E[Y(1)-Y(0)]$. Specify whether assignment changes auction prices or rival ad allocation; if so, the intervention is a market policy rather than an isolated individual impression. Holdout definitions, attribution windows, deduplication, and outcome missingness are part of the design.

\subsection*{Formal research question}

Identify campaign-specific incremental profit,
\[
\Delta\Pi=\mathbb E[M(1)-M(0)]-\mathbb E[K(1)-K(0)],
\]
where $M$ is contribution margin before advertising costs and $K$ the campaign's advertising expenditure. Both use the same fixed eligible population, horizon, and normalization: in the per-consumer version, divide each policy's total margin and total expenditure by that population's size; total incremental profit multiplies the contrast by that size. Determine designs and sample sizes sufficient for a confidence interval to resolve whether $\Delta\Pi$ exceeds a declared profitability threshold. If the target is a causal exposure effect rather than eligibility, state and assess the exclusion, monotonicity, and relevance assumptions needed for an instrumental-variable interpretation. Extend the target to total incremental effects of a multichannel campaign, allowing channel substitution and a declared interference structure.

\subsection*{Interpretation and scope}

A platform's attribution rule assigns credit to observed exposures but does not automatically estimate purchases that would disappear without advertising. Randomized holdouts identify the specified eligibility policy even when most eligible consumers never see an impression. Dividing an eligibility effect by the exposure rate identifies a local exposure effect only under additional instrumental-variable assumptions and an appropriate exposure model. Conditioning on consumers who were exposed destroys the original randomization. Ghost-ad designs can preserve opportunity-to-expose comparisons, but their auction implementation and exclusions must be explained. Geo experiments allow broader spillovers at the cost of fewer independent units; conventional individual standard errors are then inappropriate. Low conversion rates and large purchase variance can require substantial samples, so the atlas's easy label is not guaranteed. A useful extension would provide advertiser-controlled, reproducible holdouts measuring cross-platform outcomes and assess heterogeneity without selecting profitable subgroups after observing results. Generalization remains campaign-, audience-, and horizon-specific.

\subsection*{Source audit and status}

[S1] is a paid-search field experiment; [S2] is the digital-advertising measurement chapter, not a different two-author Lewis--Rao precision paper; [S3] compares observational methods against Facebook experiments. The original three citations are verified with distinct roles. Many campaign-specific causal effects are already measured; the editorial frontier is scalable, precise, multichannel measurement under targeting and interference, not an unresolved question whether randomized advertising experiments are possible.

\subsection*{References}

\begin{enumerate}[label={[S\arabic*]},leftmargin=*]

\item Thomas Blake, Chris Nosko, and Steven Tadelis (2015). \emph{Consumer Heterogeneity and Paid Search Effectiveness: A Large-Scale Field Experiment}. Econometrica 83(1), 155--174. \url{https://doi.org/10.3982/ECTA12423}\par \textit{Source role:} Randomized paid-search evidence. \textit{Locator:} Abstract and article metadata.

\item Randall Lewis, Justin M. Rao, and David H. Reiley (2015). \emph{Measuring the Effects of Advertising: The Digital Frontier}. Economic Analysis of the Digital Economy, Chapter 7, University of Chicago Press, 191--218. \url{https://www.nber.org/books-and-chapters/economic-analysis-digital-economy/measuring-effects-advertising-digital-frontier}\par \textit{Source role:} Measurement and experimentation review. \textit{Locator:} Abstract and article metadata.

\item Brett R. Gordon, Florian Zettelmeyer, Neha Bhargava, and Dan Chapsky (2019). \emph{A Comparison of Approaches to Advertising Measurement: Evidence from Big Field Experiments at Facebook}. Marketing Science 38(2), 193--225. \url{https://doi.org/10.1287/mksc.2018.1135}\par \textit{Source role:} Comparison of observational and randomized measurement. \textit{Locator:} Abstract and article metadata.

\end{enumerate}

\subsection*{Common conventions: Source mathematical conventions}

Unless an entry states otherwise, agent and firm sets are finite. State and action spaces are measurable, strategies and outcome maps are measurable, probability laws are specified, and expectations exist for the targets under discussion. Infinite-horizon discount factors lie in $(0,1)$ and discounted objectives are finite. Norms, tolerances and complexity input sizes must be specified for computational comparisons. Constants or primitives that appear locally are inputs, not empirical facts asserted by this catalogue.

\textbf{Identification.} A statistical model $m\in\mathcal M$ implies an observable law $P_m$ and target $\theta(m)$. At law $P$, its identified set is
\[
\Theta(P;\mathcal M)=\{\theta(m):m\in\mathcal M,\ P_m=P\}.
\]
Point identification holds when this set is a singleton, or equivalently when $P_{m_1}=P_{m_2}$ implies $\theta(m_1)=\theta(m_2)$ for all compatible models. Sharp bounds mean bounds attained, or approached when the boundary is not attained, by compatible models. Writing this set is a definition, not a solution: the substantive task is to establish informative restrictions and derive or estimate the set. An unrestricted-confounding model may give only trivial bounds.

\textbf{Inference.} Identification, estimability and valid uncertainty quantification are separate questions. Fix $\alpha\in(0,1)$. A confidence procedure $C_n$ over a declared law class $\mathcal P$ has asymptotic uniform coverage at level $1-\alpha$ when
\[
\liminf_{n\to\infty}\inf_{P\in\mathcal P}\Pr_P\{\theta(P)\in C_n\}\geq1-\alpha.
\]
For set-identified targets the coverage event must be defined separately, for example coverage of the full identified set. If an entry uses identification-indexed models instead of uniquely indexed laws, the same coverage convention is applied over those models. Pointwise limits do not establish uniform validity, and asymptotic validity does not guarantee finite-sample precision. Sample sizes, dependence structures and weak-identification sequences are part of each application.

\textbf{Counterfactuals.} $Y_i(a)$ is a potential outcome under a specified intervention, including its timing, implementation and versions. It is not $Y_i$ conditional on voluntarily choosing $a$. A full intervention vector permits interference; shorthand scalar treatment notation does not silently impose no interference. Assignment, selection, attrition, anticipation and measurement must be specified. A claim about an instrument's local effect is not automatically a population policy effect.

\textbf{Economic equilibrium.} A counterfactual model includes best responses, constraints, feasibility, market clearing, state transitions and expectations consistent with the stated information. When several equilibria exist, either a declared selector is an input or the target is set-valued. Comparisons across policy regimes must use the stated selection convention. Reduced-form relationships are not presumed invariant after an equilibrium-changing intervention.

\textbf{Optimization and welfare.} Optimization is over the stated feasible set. If attainment is not established, $\sup$ or $\inf$ and $\varepsilon$-optimal choices replace an asserted maximizer or minimizer. For example, with value $V=\sup_{a\in\mathcal A}W(a)<\infty$, an $\varepsilon$-optimal set is $\{a\in\mathcal A:W(a)\geq V-\varepsilon\}$ for $\varepsilon>0$. Benefits, costs, risk and health or environmental outcomes can be aggregated only using declared units and conversion weights. Interpersonal utility weights, the treatment of fiscal transfers and foreign profits, discounting, and the behavioral welfare criterion are explicit inputs. A comparative-static derivative additionally requires existence and appropriate differentiability of the selected counterfactual.

\textbf{Sources.} Local labels [S1], [S2], and so forth restart in every question. Locators identify the relevant abstract, model, section or result at the level actually checked; a bibliographic or abstract-level verification is not represented as inspection of an entire proof. Working papers are distinguished from later publications and changing versions. Source summaries are paraphrased. All URL checks and editorial formulations refer to the audit date stated above.
% END ECONOMICS STATEMENT
\end{document}
